% M.Carlsen - W.Anand, World Chess Championship: your comments. Format: games/README.md.
@ title | M.Carlsen - W.Anand, World Chess Championship

@ 6...O-O
\mythoughts{The position resembles the Italian or the Ruy Lopez, where White often reroutes
the b1-knight to f5 via g3 or e3 and builds a center with c3 and d4.
Compared with those standard structures, Black's dark-squared bishop is on c5
rather than e7. I would still choose 7.c3. If
7...a6 8.\B a4 b5 9.\B c2, then once I establish the center, Black's c7-pawn
can become a target.}

\myreflection{My move is playable, but the claim that ``Black's c7-pawn can become a target'' is wrong:
the bishop can simply protect it from b6.}

@ 7...bxc6
\mythoughts{Candidate moves are \N bd2, c3, and h3. 8.h3 is less relevant,
because Black is unlikely to play \B g4 and exchange that knight, as it
would help White open the g-file. I would choose 8.\N bd2.}

\myreflection{My move is not good.
True, Black is unlikely to play \B g4, but he can play \N g4,
so h3 does have its merits.
Also, 8.c3 \N g4 9.d4 is playable.}

@ 8...Re8
\mythoughts{The usual plan for building the center is \N bd2, c3, \Q c2, and then d4.
I still prefer 9.\N bd2.}

@ 9...Nd7 | board
\mythoughts{Black wants to play \N f8-e6 to stop me from establishing a pawn center.
If I fail to set up that center, I do not see how to improve my position.
So I need to hurry with 10.c3.}

\myreflection{My move is playable. Once I can match the position to something in my knowledge base,
I can play a series of moves without calculating each one.\index{Practical play!Pattern recognition}

Magnus Carlsen thinks only when he cannot find a similar position;
otherwise he simply reuses an existing solution.}

@ 10...Bb6
\mythoughts{I want to place my dark-squared bishop on b2, as that seems to be its best square.
I also want to play g3, \N h4, f4, \N f5 etc.
Black must watch out for sacrifices on e5 and pressure against g7.

After 11.a4 a5, I don't get anything since I have no access to b5.

I want to play 11.b3.}

\myreflection{My claim that ``after 11.a4 a5, I don't get anything since I have no access to b5''
is wrong. After the exchange on b6, Black can no longer recapture with the a-pawn,
which opens the a-file. Such an exchange therefore favors White.

My move is playable.
However, the plan of placing the bishop on b2 is not realistic: White
must reckon with Black's central strike d5:\index{Strategy!Central counterstrike}
11.b3 d5 12.exd5 cxd5 13.\N cxe5 \N xe5 14.\N xe5 \R xe5 15.\R xe5 \B xf2+ 16.\K h2 \B d4 17.\B f4 \B xe5 18.\B xe5
The line is almost forced, and I had not calculated it at all.}

@ 11...a5 | board
\mythoughts{Playing for an attack on the kingside is a typical
motif for Ruy Lopez structures. One idea would be
g3, h4, \K g2, \R h1. However, I must account for
Black can strike in the center with d5, e4, \N e5 etc.
The attack is stopped.

The other route is \N g5, h4, \R e3. To make this work,
I first need to remove the bishop on b6. Therefore 12.\N xb6.}

@ 12...cxb6 | board
\mythoughts{I do not like attacking with f4 yet because my rooks are not ready.
Black can defend f7 without difficulty.

I would like to attack with \N g5. In 13.\N g5 h6 14.\N xf7 \K xf7 15.\Q h5+ g6, I
don't see how to proceed. 13.\N g5 h6 14.\N xf7 \K xf7 15.\B xh6 gxh6 16.\Q h5+ \K e7 17.\R e3 \N f6
I still do not see a clear continuation. My pieces are not fully developed, and sacrificing
two pieces is too early.

13.h4 looks better. Black would prefer not to play h6 first because it weakens the kingside.
13.h4 \N f6 14.\B g5 the pin is unpleasant for Black.\index{Pin}
13.h4 d5 is unpleasant for White. White must wait for Black to play dxe4. Black
can then move his knight to e6 targeting d4. c3 would not help because it weakens the d3 pawn
permanently.

Clearly I need another idea. I wanted to play d4 anyway, so what about doing it now?
13.d4 exd4 14.\N xd4 hitting c6 and f5. It looks good for me.
I also intend to play 14.dxe5, 15.\B e3 targeting the b6 pawn.}

\myreflection{The variations above are based on a single idea and should be easy to defend.
In practice, I should not spend much time calculating such lines, because it is
unlikely that I will get much from them.

I was mixing ideas from the game Carlsen vs. Wang Hao, Tata Steel 2011.}

\diagram{2rqkb1r/1p1n1pp1/p1n1p1p1/2ppP3/3P4/2P1BN2/PP2BPPP/R2Q1RK1 w k - 0 13}

\myreflection{There, the center is completely closed, so the pawn thrusts g3 and h4 are effective
because Black cannot cover the weakness with his rooks.

In the actual game, even if I open lines successfully, Black can simply place
his rooks on the 7th rank and my attack is over.

A more subtle approach is to focus on the opponent's weaknesses:
tie him down to defending one weakness, then strike somewhere else where the defense
is thin. Alekhine won many games with this idea.\index{Strategy!Principle of two weaknesses}\index{Strategy!Tying down}

As for the concrete move, after 13.h4 c5 Black is better, because White
can never set up the center. The lesson here, clichéd as it is: pay attention
to a central counterstrike whenever you plan to play on the flank.\index{Strategy!Central counterstrike}}

@ 13...Qc7 | board
\mythoughts{Black's last move protects the weak pawns on the queenside.

My c1 bishop is not active yet, and I am not sure where to place it. Both a3 and e3 are
reasonable squares for it.
My a1 rook is the worst piece and therefore I want to improve it.\index{Strategy!Improve the worst piece}
Since I do not yet know where to develop the c1 bishop, I do not want to keep
the rook on the first rank; it should wait until the bishop is developed.

Developing the rook via a3 looks good. I can then place it on d3 and target d6. My
bishop can either go to a3, e3 or even g5 depending on what happens next.

I do not see any immediate threat from Black. 14...d5 drops a pawn after
15.dxe5 \N xe5 16.exd5 cxd5 17.\Q xd5

So I decide to play 14.\R a3.}

\myreflection{Yes, I found the same move as Carlsen!

By the way, the engine recommends 14.c4, which is the same idea as my next move.}

@ 14...Nf8
\mythoughts{I would love to place my knight on f5. If Black's bishop takes it,
I recapture with the e-pawn, and f6 can become a serious threat.

That idea does not work because Black can take on d4 with exd4.
I have a weakness on the first rank and could lose a bishop.

This means the center must be stable before I take any action\index{Strategy!Central counterstrike}
on the kingside. c4 and then d5 is natural for me.

Another idea is \R d3. After exchanging the e-pawns,
the heavy pieces will be all exchanged. Do I have advantage with
the remaining knight and bishop endgame? Black can place his
bishop on e6, denying the c4 square for my knight. We both have
a good bishop. I do not see how I can break through, since my king has no route
through the light squares; they are all covered by
Black's bishop.

I decide to play 15.c4.}

\myreflection{15.\R d3 is playable, but White should then close the center with d5.

After 15.\R d3 \N g6 16.dxe5?! dxe5, Black need not exchange pieces, for the same reason: White's
heavy pieces have no targets. It is therefore unsound to speculate about an endgame
that would not actually arise here.}

@ 15...dxe5
\mythoughts{I thought I had to place my rook on d3, otherwise Black would occupy
the d-file. After deeper consideration, I no longer think it is mandatory:
my first and second ranks are covered. Black's rooks have no squares to invade.

If I still want to attack on the kingside, I need to move my f3-knight
so that the a3-rook can swing over.

16.\N h4 \R d8 17.\Q h5 is natural. All my pieces can join the
attack soon.

There is a clear lesson here: occupying an open file does not always bring
an advantage when the rooks have no targets.\index{Strategy!Entry points for rooks}}

@ 17...f6
\mythoughts{I realize 18.\N f5? g6 I will lose material.
18.\R g3 is natural. Next I can bring my knight to f5. Black still needs to develop the
bishop. I feel I have some advantage now.}

\myreflection{The statement above is wrong:
18.\N f5 g6? 19.\R g3 \K h8 20.\Q h4 \B xf5 21.\Q xf6+ \Q g7 22.\Q xg7+ \K xg7 23.exf5.

My intended move 18.\R g3 is still good, however, since I can follow up with 19.\N f5.}

@ 18...Be6
Incidentally, this is a poor move: the e6 square belongs to the knight, not the bishop.

\mythoughts{I don't need to worry about 19...\B xf5. In the endgame
this bishop is the most important defender. When it gets traded,
I will have a winning endgame.

I still want to play 19.\R g3.}

@ 19...Ng6 | board
\mythoughts{I still need more pieces for the attack. Concrete
calculation is not yet necessary. This saves me time and energy.

Bringing my e1-rook to the g-file does not help much. It costs tempi and
allows Black to counterattack with \R d1+.

Even without that counterattack, I can have at most four attackers
against the g7 pawn (two rooks on g-file, one f5 knight and queen or bishop on h6).
Black already has four defenders.\index{Strategy!Attackers against defenders} So this brute-force attack does not work
out.

The same argument also applies to my c1 bishop, since it can only attack g7 as
its target.

Through this process of elimination,\index{Practical play!Process of elimination} I conclude I must advance the h-pawn. Once it reaches h6,
Black may be forced to play g6, which allows further piece sacrifice on g6.

I must move my queen to make room for the h-pawn. Placing the queen on g4 lets her
protect the h-pawn on h4; otherwise Black can take it by first exchanging the knight
on f5 and then take on h4.

I decide to play 20.\Q g4.}

\myreflection{I like that I am making logical decisions, even if the conclusion is not entirely correct.
After 20.\Q g4 \K h8 21.h4 \B xf5 I should recapture with the queen rather than the e-pawn, because otherwise Black
can play 22...\R d4.

After 20.\Q g4 \K h8 21.h4 \B xf5 22.\Q xf5 \N xh4 23.\Q h5 g5 24.\B xg5 fxg5 25.\R xg5 \N xg2 26.\R xg2 \R g8
the game is equal.
So here it does not matter whether I place the queen on g4 or f3, as both lead
to the same variation.}

@ 20...Bxf5
\mythoughts{This is indeed a strange move. Black has already lost a tempo
with the bishop, and exchanging it for my knight is also a bad idea,
since I have established that the endgame favors me.

I decide to play 21.exf5. It drives the knight away, and a pawn
on f5 will also help me after I play h6 and Black responds with g6.}

@ 21...Nf4
\mythoughts{I decide to play 22.\Q g4. I want my queen on the kingside for the attack.
I don't see any other squares for her.}

\myreflection{The queen on h5 is actually well placed: as we will see, it controls
both the e8 and the d1 squares. I failed to reason soundly here.
Even though the queen needed a better square, g4 is not it, since
I had already established that I do not have enough attackers against g7.

The text move is quite logical: remove the best defender.\index{Strategy!Remove the best defender}
Besides, my bishop has not moved at all, and the knight on f4 cost Black
a lot of tempi, so I have a clear lead in development.}

@ 22...exf4
\mythoughts{I have calculated 23.\R c3 \R d5 24.\Q f3 \Q b7 25.\R e6 \R c8 26.\R exc6 \R xc6 27.\Q d5+ winning material.
I don't like 23.\R c3 c5 because my rook is not doing
a lot on c3.

I don't like 23.\R g4 at all because my queen is stuck on h5.

I do not like 23.\R f3 either. I need more tempi to win the f4-pawn.
My pieces are cramped. It does not look
like a good deal.

I don't like 23.\R g6 because the move is too committal: my queen and rook
are stuck on the kingside and I don't see how to proceed.

I must play 23.\R c3; it is the least painful option.}

\myreflection{23.\R c3 \R d5 24.\R e8+ \R xe8 25.\Q e8\# mate! As we will see,
this is not the only time I miss this simple tactic.\index{Back-rank mate}

23.\R f3}

@ 23...c5 | board
\mythoughts{My queen on h5 doesn't do much and therefore I want to improve it.
24.\Q f3 is a natural move. 25.\R c4 will be the next move,
winning the f4 pawn.}

\myreflection{The f4-pawn is not important at all. The queen on f3 is worse than on h5,
as we have already established.}

@ 24...Rab8
\mythoughts{25.\Q e2 threatens 26.\R e7, which is a natural follow-up.}

@ 25...Qd7 | board
\mythoughts{Black threatens to play 26...\Q d1+ 27.\Q xd1 \R xd1+ 28.\K h2 \R d2 Black
has an advantage in the endgame because the rook on d1 is more active.

I do not see how 26.\K h2 \Q d1 27.\Q xd1 \R xd1 helps me. My f2-pawn is too weak.
I want to play 26.\R ce4 \Q d1+ 27.\Q xd1 \R xd1+ 28.\K h2 \R d2 29.\R e2 \R xe2 30.\R xe2.
My king arrives in time to protect the f5-pawn and take the f4-pawn. Black's rook
cannot stay on d4 because I can play b3 and then c3, kick it away.}

\myreflection{``I do not see how 26.\K h2 \Q d1 27.\Q xd1 \R xd1 helps me.'' It is
painful to miss winning material with 27.\R e8+ \R xe8 28.\Q xd1.

My calculation is correct, but my evaluation of the rook endgame is also wrong.
After 26.\R ce4 \Q d1+ 27.\Q xd1 \R xd1+ 28.\K h2 \R d2 29.\R e2 \R xe2 30.\R xe2 b5, the endgame is equal.}

@ 26...Rf8
\mythoughts{I must keep queens on the board. If Black manages to exchange queens
on d1, his rook on d1 cuts my king off. I don't think I can win this endgame.

27.\R ce4 prevents Black from playing 27...\R be8. I threaten
28.\R e7 in the next move.
27.\R ce4 \Q c7 28.\R e7 is winning for White.

27.\R xf4 is worse since the f4-pawn is not the focus of the game.

I decide to play 27.\R ce4.}

\myreflection{I am glad to see that I recognized the f4-pawn is not the focus of the game.}

@ 27...Rb7
\mythoughts{I want to play 28.\Q e2, threatening 29.\R e7 \Q c8 30.\R xb7 \Q xb7 31.\R e7 \Q c8 32.\Q g4 \R f7 33.\R xf7 \K xf7 34.\Q xf4.
29.\R e7 \Q c8 30.\Q g4 is even better because Black cannot answer with 32...\R f7.}

@ 28...b5
\mythoughts{I do not see what Black is up to. I want to carry out my planned
29.\R e7.}

\myreflection{I could not calculate the whole line:
29.\R e7 \Q d6 30.f3! \R xe7 31.\R xe7 bxa4 32.\Q e4 \Q b8 33.\Q xa4 \Q d6 34.\R e4 h6 35.\Q xa5 \K h7
Here f3 is an important prophylactic move.\index{Strategy!Prophylaxis}}

@ 29...bxa4 | board
\mythoughts{I will take the pawn with 30.\R xa4. Allowing Black to play ...a3
does not make sense since I do not have a mating attack. Taking with the b-pawn opens
the b-file.}

\myreflection{My move keeps the b-file closed, but misses the more important e-file:
30.\R xa4 f3 31.\Q xf3 \R b4 32.\R a6!
White still has enough advantage, but I am not sure I would find the winning line after
31...\R b4.}

@ 30...Rb4
\mythoughts{Suddenly I see 31.\R e7. Black must take the a4 pawn otherwise
32.\R xb4 axb4 33.\Q c4+ \K h8 34.\Q f7.
After 31...\Q xa4 32.\R xg7+ \K xg7 33.\Q g4+ \K h6, I don't see a checkmate.

32.\Q g4 \R f7 33.\R e8+ \R f8 34.\R 4e7 White is winning.}

@ 31...Qd6
\mythoughts{The original plan doesn't work since Black's queen covers the f8 square.
32.\Q f3 threatens 33.\R xb4 and then 34.\Q b7.}

@ end
\subsection*{My Reflection}

I am glad I was able to make moves based on logic. This way, I can arrive at
good decisions without calculating a lot.

During the game I also learned that occupying an open file is not always a good idea
when the rooks have no target.\index{Strategy!Entry points for rooks} With this in mind, I should no longer assume that
occupying an open file is automatically the right move.

Missing the back-rank tactic was painful:\index{Back-rank mate} because of that oversight, I was forced
to settle for suboptimal moves.

